\section{Discussion}
The holdout supports a positive meeting-window return, particularly relative to matched non-meeting days, but does not replicate the development finding that the premium is concentrated in anxious states. Three readings are consistent with the data. First, the premium may be roughly state-independent after 2020; the calm cell mean rose from $-6.6$ to 40.4~bps and the overall mean from 2.8 to 38.3~bps. Second, the state contrast may be present but undetectable: the calm cell has nine meetings and the minimum detectable effect for the anxious half is 40.3~bps. Third, the 2020--2022 episode, in which VIX was high and the communication regime changed, may dominate a seven-cluster sample. The data cannot separate these.

Limitations are as stated in the design: release times are assigned by rule, and for eight 2011--2012 development meetings the rule is wrong (the statement came at about 12:30~ET), which we corrected after the holdout run (Section~\ref{sec:amend}), clustering by year is a proxy for volatility episodes, and the VIX-change and press-conference tests are exploratory. The comparison with the post-publication decline in \citet{mclean2016} is loose: the holdout premium does not fall, but an announcement-window return is not a cross-sectional anomaly portfolio, and the sample is too small to measure a decline precisely.
